\documentclass{article}
\usepackage{graphicx}
\usepackage{amsfonts}
\usepackage{amsmath}

\title{Performance éngergétique du code}
\author{Gabriel xxx, Samuel MICHEL-CHAGNOT}
\date{17 Avril 2026}
\begin{document}

\maketitle

\section{Introduction}
\indent Dans la continuitée du travail fait dans le rapport précédent, on se propose ici d'étudier l'impact énergétique provoqué par les méthodes ou approches de codage.
Nous allons donc procéder à l'éxecution d'un même algorithme en utilisant divers méthodes et approches : langage compilé à la volée(Python), langage compilés (C++).

De la sorte, nous utiliserons le programme \texttt{s-tui} afin de faire des relevés énergétiques de l'éxecution, la compilation de code.
Notre objet d'étude sera le calcul de l'ensemble de Mandelbrot.
Enfin, le processeur utilisé est un Intel Xeon W1270 dont le fréquence d'horloge est de 3.40 GHz.

\section{Langage compilé à la volée}
Le sujet nous propose d'utiliser directement un script Python calculant une image représentant les fracatales de Mandelbrot.

Pour le calcul de la quantitée de CO2e du programme, nous nous basons sur les présentes dans le site elecctricitymaps.com. Nous avons ainsi défini que $1 kWh$ représente $17gCO2e$.
Enfin, pour chaque exécution, nous relevons constamment les données suivantes :
\begin{itemize}
	\item Durée d'éxecution (s)
	\item Puissance active moyenne (W)
	\item Puissance moyenne au repos (W)
	\item Energie absorbée par  la tâche (mWh)
	\item Quantité d'émission de CO2e (gCO2e)
\end{itemize}

\end{document}




















